\documentclass[11pt]{article}
\usepackage[margin=1in]{geometry}
\usepackage[T1]{fontenc}
\usepackage{lmodern,mathptmx}
\usepackage{xcolor,booktabs,tabularx,titlesec}
\titleformat{\section}{\large\bfseries}{}{0pt}{}
\titleformat{\subsection}{\normalsize\bfseries}{}{0pt}{}
\setlength{\parindent}{0pt}\setlength{\parskip}{0.6em}
\newcommand{\reviewer}[1]{{\color{blue!55!black}\itshape #1\par}}
\begin{document}
\begin{center}{\Large\bfseries Response to Reviewers}\\[0.4em]
Cost-Aware Visual Router (CAVR): Dynamic Query Routing for Cost-Efficient Multimodal RAG \textcolor{red}{[AUTHOR ACTION: update if the title changes]}\end{center}
\textcolor{red}{\small DRAFT. Reviewer comments below are paraphrased from the revision plan; replace each with the reviewer's verbatim text before sending. Items marked [AUTHOR ACTION] must be completed or removed. Tables and figures are named, not numbered; insert the final numbers from the manuscript.}

We thank the editor and the three reviewers for a careful and demanding review. The reviewers found inconsistent numbers, a cost model that contradicted our own measurements, and released code that did not match the paper. They were right on each point, and we therefore rebuilt the experimental pipeline from scratch and re-ran every experiment.

Every answer path (text, visual and hybrid) is now run on every one of the 8,150 DocVQA and InfographicVQA validation questions, with real BM25 and ColPali retrieval and one local generator (Qwen2.5-VL-7B-Instruct). Answer quality (ANLS and exact match), GPU-seconds, energy and latency are measured per query on one GeForce RTX 4090. Route labels now come from measured outcomes rather than a keyword rule, and routers are evaluated on held-out pages of five page-grouped splits. As the editor requested, the paper is now framed as a GPU resource-scheduling problem, with new experiments on serving capacity under load, the GPUs needed at a latency target and visual-index scaling.

Several claims of the submitted version are withdrawn, and we list them first so that nothing is hidden behind the new results. The title, abstract, discussion and conclusion were rewritten around what the new measurements support. All tables and figures are generated automatically from the released replication code and cached outputs, so every number can be checked. \textcolor{red}{[AUTHOR ACTION: tag the release, archive it with a DOI and cite it.]}

\section*{Withdrawn or corrected claims}
\begin{tabularx}{\textwidth}{>{\raggedright\arraybackslash}p{0.27\textwidth}l>{\raggedright\arraybackslash}X}
\toprule \textbf{Submitted claim} & \textbf{Status} & \textbf{What the revised paper reports} \\ \midrule
86.2\% routing accuracy (231/268) & Withdrawn & The keyword labels behind it agree with outcome-based labels at Cohen's $\kappa$ = 0.01, so routing accuracy is no longer reported. The submitted confusion matrix gave 229/268 = 85.4\%, which was also misreported. \\
\$0.00019 per query; 21.5\% cost saving & Withdrawn & Both came from word count $\times$ GPT-3.5 prices and router latency $\times$ an A100 rate, with no retrieval or generation run. Costs are now measured. \\
Text path about 100$\times$ cheaper than visual & Withdrawn & With a shared generator, text costs 0.43 and visual 0.62 GPU-s per query. \\
92.6\% saving on the natural distribution & Withdrawn & Never measured. All results now use the full validation sets at their natural question mix. \\
Out-of-distribution generalisation on InfoVQA & Withdrawn & InfoVQA was also used for training; both datasets are now in-distribution and split by page. \textcolor{red}{[AUTHOR ACTION: confirm, or add a held-out benchmark.]} \\
Training details (AdamW, loss weights $\alpha$/$\beta$/$\gamma$, temperature 0.7) & Corrected & The released code trained a plain classifier. The revised routers and their hyperparameters are documented and released. \\
Results without released code (temperature and $\alpha$ sweeps, bootstrap intervals, imbalanced baseline) & Removed & Replaced by experiments generated by the released code. \\
``Production-ready'' system & Reworded & ``Deployable system design'', supported by the serving experiments. \\
\bottomrule\end{tabularx}

\section*{Editor}
\subsection*{Comment E.1: Supercomputing, real-time and parallel justification}
\reviewer{Justify, with evidence, the need for supercomputing, real-time and parallel or distributed computing in the Abstract, the Introduction and throughout the paper.}
\textbf{Response.} The revised paper frames CAVR as a scheduler for GPU resources in multimodal RAG serving, and the Introduction opens with a new subsection, ``Why routing in multimodal RAG is a supercomputing problem''. Generation dominates per-query compute, so sending every query through the most expensive path inflates the number of accelerators needed to meet a latency target. In a Poisson load test with dynamic batching, simulated from measured service times, serving 100 queries/s at a p95 latency of 2 s needs 106 GPUs for always-hybrid and 72 for the content-aware router CAVR-CE (32\% fewer; mean over five splits).

Late-interaction visual retrieval stores 1,031 vectors per page, so the index grows with the corpus: at 100k pages it needs 24.6 GiB, more than our 24 GiB GPU holds, and a million pages (about 246 GiB) would exceed any single accelerator. Such an index must be streamed from host memory (1,044 ms per query on our hardware), pooled (recall@1 drops from 0.577 to 0.525) or sharded. We state that routing saves generation compute rather than index memory.

New material: a hardware table, energy per query from NVML, router throughput up to batch 1,024, serving capacity and GPUs needed per policy, visual-index latency from 1k to 100k pages, and serial versus speculative escalation. The method figure was redrawn as the pipeline actually built, including its deployment on a GPU pool. We are explicit about scale: all measurements use one GPU, capacity beyond it is simulated from measured service times, multi-GPU sharding is emulated, and larger indexes are built from replicated page embeddings. \textcolor{red}{[AUTHOR ACTION: write and verify the related-work citations on inference serving and late-interaction retrieval.]}

\subsection*{Comment E.2: Accurate reporting and limitations}
\reviewer{Report the results accurately, rewrite overstated conclusions and explain the limitations fully.}
\textbf{Response.} Every number now comes from one cached run of every path on every query, and tables and figures are generated from the results files. The withdrawn claims are listed above. Policy results are means and standard deviations over five page-grouped splits, with McNemar and paired-bootstrap tests, and every comparison states its cost. A new Limitations subsection covers the single GPU, simulated rather than live traffic, escalation latencies computed from measured per-stage times, synthetic index scaling, emulated sharding, GPU-only energy, text-only queries, and the 35\% of queries that no path answers.

\section*{Reviewer 1}
\subsection*{Comment 1: Latency of perplexity-based escalation}
\reviewer{Escalation based on generation perplexity forces a generate-then-escalate chain, so escalation can be slower than choosing the visual path from the start.}
\textbf{Response.} We agree, and the revised Method separates three escalation designs with per-query latency and cost formulas. Signals available before generation (router confidence, ColPali score margin) re-route without a serial step: on CAVR-CE, re-routing on the ColPali margin at $\tau$ = 0.5 raises p95 latency from 1.24 s to 1.36 s, against 1.39 s for always-hybrid. Answer-level signals (token log-probabilities) require the first answer: a serial cascade has p95 1.91 s, and launching the hybrid path speculatively alongside the routed path cuts p95 to 1.42 s but costs 1.46 GPU-s per query, 45\% more than always-hybrid. These latencies are computed from measured per-stage times; two concurrent workers were not run. We recommend pre-generation escalation.

\subsection*{Comment 2: Runtime budget control}
\reviewer{The cost weights are fixed, so the cost--quality trade-off cannot be changed at run time without retraining.}
\textbf{Response.} The router now predicts per-path answer quality $\hat{q}_p$, and the decision rule $\arg\max_p$ $\hat{q}_p$ $-$ $\lambda\cdot\bar{c}_p$ takes the cost weight $\lambda$ as a run-time input. One trained router therefore covers a range of budgets: sweeping $\lambda$ from 1 to 0 moves CAVR-CE from ANLS 0.502 at 0.56 GPU-s to 0.564 at 0.84 GPU-s (mean over splits). A cost--quality frontier figure and a table of operating points with their serving capacity were added. Adapting $\lambda$ to load is described as a deployment option but not evaluated. \textcolor{red}{[AUTHOR ACTION: decide whether to add a load-adaptive $\lambda$ experiment.]}

\subsection*{Comment 3: Queries containing images}
\reviewer{The query encoder is text-only, so queries that contain an image are not supported.}
\textbf{Response.} We narrowed the scope to text queries over visual documents and state this in the title, the abstract and the Limitations. \textcolor{red}{[AUTHOR ACTION: confirm narrowing the scope, or add a proof of concept with an image embedding.]}

\section*{Reviewer 2}
\subsection*{Comment 1: Estimated versus measured costs}
\reviewer{The estimated and measured costs are inconsistent. Which of them produces the reported 21.5\% saving?}
\textbf{Response.} Neither is used any more. The 21.5\% came from list-price estimates, which contradicted the measured costs, and neither involved real retrieval or generation. Costs are now measured per query on one GeForce RTX 4090: over all queries, always-text costs 0.43, always-visual 0.62 and always-hybrid 1.01 GPU-s (144, 170 and 329 J). All savings are computed from these measurements, including the router's own overhead and, where the router reads ColPali scores, ColPali retrieval for every query.

\subsection*{Comment 2: Equivalence of the objectives}
\reviewer{Constrained cost minimisation is not equivalent to maximising the quality-to-cost ratio.}
\textbf{Response.} We agree. The ratio objective is replaced by a Lagrangian penalty: the router picks $\arg\max_p$ $\hat{q}_p$ $-$ $\lambda\cdot\bar{c}_p$. The Method now derives this as the Lagrangian relaxation of budget-constrained routing, states when it is exact (randomised policies over a finite query set, by linear-programming duality), notes that the best deterministic policy is within 1/N of that optimum, and gives an example in which the ratio rule and the constrained rule disagree. \textcolor{red}{[AUTHOR ACTION: check the derivation.]}

\subsection*{Comment 3: How labels are produced}
\reviewer{Explain how the optimal-path labels and utility targets are produced, and report the number of samples per class.}
\textbf{Response.} In the submitted version, labels came from a keyword rule applied to InfographicVQA questions; we state this plainly. These labels agree with outcome-based labels at Cohen's $\kappa$ = 0.01, i.e. no better than chance. Labels are now produced by running every path on every query, scoring each answer with ANLS and taking the path that maximises ANLS $-$ $\lambda\cdot\bar{c}_p$, with ties broken towards the cheaper path. Class counts per split and a threshold-based variant are reported in the labels table.

\subsection*{Comment 4: Escalation ablation, calibration and the threshold}
\reviewer{Provide an ablation of the escalation signals, reconcile the two reported ECE values (0.06 and 0.042), and justify the threshold $\tau$ = 0.6.}
\textbf{Response.} The ablation covers no escalation, each signal alone, all pre-generation signals and all signals, for both routers, with $\tau$ swept from 0.3 to 0.9 (escalation table and figure). We report the whole $\tau$ curve instead of a single tuned value. The conflicting ECE values are replaced by test-set measurements with 10 bins: the CAVR confidence head has ECE 0.052, and the escalation probabilities, recalibrated on the validation split by logistic regression, have ECE 0.027--0.078 depending on the signal. \textcolor{red}{[AUTHOR ACTION: add a reliability diagram if requested.]}

\subsection*{Comment 5: Minor points}
\reviewer{Table 5 columns are reversed; ``79'' versus ``+79.9 pp''; ``two orders of magnitude'' is about 50$\times$; garbled text (``volto brio''); ``production-ready'' is too strong.}
\textbf{Response.} Table 5 is removed, and targets are no longer reported. The sentence with ``79'' was removed with the withdrawn results. Router overhead is now reported as measured, 11.7 ms for CAVR and 24.8 ms for CAVR-CE per query. ``Production-ready'' is now ``deployable system design''. \textcolor{red}{[AUTHOR ACTION: remove ``two orders of magnitude'' wherever it remains, complete a full copy-edit pass, and address any further comments in Reviewer 2's attachment.]}

\section*{Reviewer 3}
\subsection*{Comment 1: Answer correctness and escalation switched off}
\reviewer{Answer correctness is never checked, and there is no comparison with escalation switched off.}
\textbf{Response.} This was the central gap and is now the core experiment. Every path is run on every query, answers are scored with ANLS and exact match, and every policy is evaluated on the answers it produces (main-results table). With $\lambda$ chosen on validation data, CAVR-CE uses 0.716 instead of 1.009 GPU-s per query (29\% less than always-hybrid); its ANLS is 0.557 against 0.566, lower on 4 of 5 splits with the 95\% interval excluding zero on 1, and the exact-match difference is not significant on 5 of 5. Escalation switched off is the ``none'' row of the escalation table.

\subsection*{Comment 2: Inconsistent numbers}
\reviewer{The paper reports 231/268 where the confusion matrix gives 229/268, and the rule-based baseline, path counts, training curve and calibration differ between sections.}
\textbf{Response.} All of these came from results that could not be reproduced, and they have been removed. Every number in the revision comes from one cached run, and tables and figures are generated from the results files. Policy results are mean $\pm$ standard deviation over five splits, and figures that show a single split are labelled ``seed 1''.

\subsection*{Comment 3: Cost ratios}
\reviewer{Visual retrieval is about 33$\times$ more expensive than text, not 100$\times$, and Table 8 even shows visual cheaper than text.}
\textbf{Response.} The reviewer is right that the cost claims were inconsistent. With measured costs over all queries, text and visual cost 0.43 and 0.62 GPU-s per query, and the paper no longer claims a large text--visual cost gap. The real cost differences are between one context and two (hybrid costs 1.6$\times$ the visual path) and the size of the visual index at scale.

\subsection*{Comment 4: Objective and labels}
\reviewer{The objective equivalence is incorrect, and the labels and utility targets are unclear.}
\textbf{Response.} Please see our responses to Reviewer 2, Comment 2 (Lagrangian formulation) and Comment 3 (outcome-based labels).

\subsection*{Comment 5: Repository does not match the paper}
\reviewer{The released repository uses InfoVQA only, keyword labels, a classification-only model and no real visual retrieval.}
\textbf{Response.} The reviewer's audit was correct on every point. The revision uses a new pipeline with DocVQA and InfographicVQA, real BM25 and ColPali retrieval, generation for every path, outcome-based labels, multi-task routers and released scripts for every table and figure. The method figure now shows the pipeline exactly as built. \textcolor{red}{[AUTHOR ACTION: tag the release and archive it with a DOI.]}

\subsection*{Comment 6: References}
\reviewer{The arXiv identifiers of [9] and [29] are wrong, and several references are incorrect or duplicated.}
\textbf{Response.} \textcolor{red}{[AUTHOR ACTION: audit every reference against its DOI or arXiv page; fix references 9 and 29, remove duplicates, correct authors and venues, and remove any entry that cannot be found. Rebuild Table 1 afterwards, because RouteLLM and FrugalGPT do model cost.]}

We hope that the revised manuscript addresses the editor's and reviewers' concerns, and we thank them again for comments that substantially improved the work.
\end{document}
